\documentclass[12pt,a4paper]{article}

\usepackage[a4paper,margin=17mm]{geometry}
\usepackage[T2A]{fontenc}
\usepackage[utf8]{inputenc}
\usepackage[russian]{babel}
\usepackage{amsmath,amssymb,mathtools}
\usepackage{array,booktabs,longtable}
\usepackage{xcolor}
\usepackage{hyperref}
\usepackage{tikz}

\hypersetup{
  colorlinks=true,
  linkcolor=blue!55!black,
  urlcolor=blue!55!black
}

\setlength{\parindent}{0pt}
\setlength{\parskip}{0.55em}
\renewcommand{\arraystretch}{1.25}

% Сообщение Софии по итогам ревью:
% София, в первой части в пункте а возникла ошибка при применении формулы приведения.
% Во втором упражнении длина гипотенузы найдена верно, но синус и косинус внешнего угла не вычислены.
% В пятом упражнении в пункте г выражение упрощено верно, но при записи области допустимых значений потерян знак.

\begin{document}

\begin{titlepage}
\centering
\vspace*{0.25\textheight}

{\LARGE\bfseries Ревью домашней работы\par}

\vspace{2.2cm}

{\large Автор: Лёвин Артём Александрович\par}

\vspace{0.9cm}

{\large 20 сентября 2026 года\par}

\vfill

\begin{tikzpicture}[x=1cm,y=1cm]
  \draw[blue!45!black, line width=0.8pt]
    (0,0) .. controls (2,0.45) and (4,-0.45) .. (6,0)
          .. controls (8,0.45) and (10,-0.45) .. (12,0);
\end{tikzpicture}

\vspace*{0.8cm}
\end{titlepage}

\newpage
\section*{Сводная таблица}

{\small
\begin{longtable}{
  >{\raggedright\arraybackslash}p{0.07\linewidth}
  >{\raggedright\arraybackslash}p{0.14\linewidth}
  >{\raggedright\arraybackslash}p{0.25\linewidth}
  >{\raggedright\arraybackslash}p{0.15\linewidth}
  >{\raggedright\arraybackslash}p{0.15\linewidth}}
\toprule
\textbf{№} & \textbf{Тема} & \textbf{Суть ошибки} & \textbf{Ваш ответ} & \textbf{Правильный ответ} \\
\midrule
\endfirsthead
\toprule
\textbf{№} & \textbf{Тема} & \textbf{Суть ошибки} & \textbf{Ваш ответ} & \textbf{Правильный ответ} \\
\midrule
\endhead
\hyperref[task:1a]{1а} &
Формулы приведения &
В формуле для синуса дополнительного угла оставлен косинус исходного угла, а не второго слагаемого. &
$\cos 30^\circ$ &
$\dfrac12$ \\
\addlinespace
\hyperref[task:2]{2} &
Внешний угол прямоугольного треугольника &
Гипотенуза найдена, но требуемые синус и косинус внешнего угла не вычислены. &
не указан &
$\sin\angle FAC=\dfrac35$, \newline
$\cos\angle FAC=-\dfrac45$ \\
\addlinespace
\hyperref[task:5d]{5г} &
ОДЗ рационального выражения &
Само выражение упрощено верно, но при решении условия знаменателя потерян знак минус, поэтому исключено неверное значение. &
$-1,\ x\ne-\dfrac35$ &
$-1,\ x\ne\dfrac35$ \\
\bottomrule
\end{longtable}
}

\newpage
\vspace*{\fill}
\begin{center}
{\LARGE\bfseries Разбор ошибок}
\end{center}
\vspace*{\fill}

\newpage
\phantomsection\label{task:1a}
\section*{Задание 1а}

\subsection*{Текст задания}
Вычислите:
\[
\sin 30^\circ.
\]

\subsection*{Ваше решение}
Записано:
\[
\sin 30^\circ=\sin(90^\circ-60^\circ)=\cos 30^\circ.
\]

\subsection*{Ваш ответ}
\[
\cos 30^\circ.
\]

\subsection*{Ошибка}
\textcolor{red}{Ошибка:} неверно применена формула приведения для синуса дополнительного угла.

\subsection*{Где именно возникает ошибка}
Последний правильный шаг:
\[
\sin 30^\circ=\sin(90^\circ-60^\circ).
\]

Первый неправильный шаг:
\[
\sin(90^\circ-60^\circ)=\cos 30^\circ.
\]

\subsection*{Почему этот шаг неверен}
Для угла вида $90^\circ-\alpha$ действует правило
\[
\sin(90^\circ-\alpha)=\cos\alpha.
\]
Здесь роль $\alpha$ играет именно $60^\circ$, потому что внутри скобок стоит $90^\circ-60^\circ$.
Поэтому после перехода от синуса к косинусу должен остаться угол $60^\circ$, а не $30^\circ$.

Иными словами, число $30^\circ$ здесь является уже результатом вычитания, а в формуле приведения после смены функции сохраняется угол $\alpha$, то есть второе слагаемое исходной записи.

\subsection*{Ключевая идея}
\textcolor{blue!60!black}{Ключевая идея:} сначала выделите стандартную форму $90^\circ-\alpha$, затем меняйте синус на косинус, сохраняя именно угол $\alpha$.

\subsection*{Исправление от последнего правильного шага}
Продолжаем с верной записи:
\[
\sin 30^\circ
=\sin(90^\circ-60^\circ)
=\cos 60^\circ
=\frac12.
\]

\subsection*{Полное правильное решение}
Представим угол $30^\circ$ как разность:
\[
30^\circ=90^\circ-60^\circ.
\]
Тогда
\[
\sin 30^\circ=\sin(90^\circ-60^\circ).
\]
По формуле приведения
\[
\sin(90^\circ-\alpha)=\cos\alpha,
\]
поэтому
\[
\sin(90^\circ-60^\circ)=\cos60^\circ.
\]
Табличное значение:
\[
\cos60^\circ=\frac12.
\]
Следовательно,
\[
\boxed{\sin30^\circ=\frac12}.
\]

\subsection*{Проверка}
Табличное значение синуса угла $30^\circ$ равно $\frac12$, поэтому полученный результат совпадает с известным значением.

\subsection*{Итог}
\textcolor{teal!60!black}{Итог:}
\[
\sin30^\circ=\frac12.
\]

\subsection*{Задача на закрепление}
Вычислите $\sin60^\circ$, представив угол как $90^\circ-30^\circ$.

\newpage
\phantomsection\label{task:2}
\section*{Задание 2}

\subsection*{Текст задания}
В прямоугольном треугольнике $ABC$
\[
\angle B=90^\circ,\qquad AB=8,\qquad BC=6.
\]
Луч $AF$ является продолжением луча $AB$ в противоположную сторону.

Найдите
\[
\sin\angle FAC
\quad\text{и}\quad
\cos\angle FAC.
\]
В решении обязательно укажите длину гипотенузы $AC$.

\subsection*{Ваше решение}
На рисунке отмечены катеты длиной $8$ и $6$ и записано:
\[
AC=10.
\]
Значения $\sin\angle FAC$ и $\cos\angle FAC$ далее не указаны.

\subsection*{Ваш ответ}
Не указан.

\subsection*{Ошибка}
\textcolor{red}{Ошибка:} решение остановлено после нахождения гипотенузы; обязательные значения тригонометрических функций внешнего угла не получены.

\subsection*{Где именно возникает ошибка}
Последний корректный результат:
\[
AC=10.
\]
Это действительно длина гипотенузы треугольника со сторонами $6$, $8$ и $10$.

Неполный шаг возникает сразу после этого: нужно перейти от внутреннего угла $\angle BAC$ к внешнему углу $\angle FAC$, однако этот переход и вычисления отсутствуют.

\subsection*{Почему решение неполное}
По условию требуется найти не только гипотенузу, но и две тригонометрические функции внешнего угла.

Луч $AF$ направлен противоположно лучу $AB$. Поэтому углы $\angle FAC$ и $\angle BAC$ являются смежными:
\[
\angle FAC=180^\circ-\angle BAC.
\]
Для смежных углов синус сохраняет знак и значение, а косинус меняет знак:
\[
\sin(180^\circ-\alpha)=\sin\alpha,
\qquad
\cos(180^\circ-\alpha)=-\cos\alpha.
\]
Именно поэтому нельзя ограничиться вычислением тригонометрических функций внутреннего угла: для косинуса внешнего угла обязательно появляется знак минус.

\subsection*{Ключевая идея}
\textcolor{blue!60!black}{Ключевая идея:} сначала найдите гипотенузу и синус с косинусом внутреннего угла при вершине $A$, затем учтите, что требуемый угол является смежным с ним.

\subsection*{Исправление от последнего правильного шага}
Имеем
\[
AC=10.
\]
Во внутреннем прямоугольном треугольнике:
\[
\sin\angle BAC=\frac{BC}{AC}=\frac6{10}=\frac35,
\]
\[
\cos\angle BAC=\frac{AB}{AC}=\frac8{10}=\frac45.
\]
Так как
\[
\angle FAC=180^\circ-\angle BAC,
\]
то
\[
\sin\angle FAC=\frac35,
\qquad
\cos\angle FAC=-\frac45.
\]

\subsection*{Полное правильное решение}
Треугольник $ABC$ прямоугольный, поэтому по теореме Пифагора:
\[
AC^2=AB^2+BC^2.
\]
Подставляем длины:
\[
AC^2=8^2+6^2=64+36=100.
\]
Длина стороны положительна, значит
\[
AC=10.
\]

Теперь рассматриваем внутренний угол $\angle BAC$.

Противолежащий ему катет равен $BC=6$, а гипотенуза равна $AC=10$. Поэтому
\[
\sin\angle BAC=\frac{BC}{AC}=\frac6{10}=\frac35.
\]

Прилежащий к этому углу катет равен $AB=8$. Поэтому
\[
\cos\angle BAC=\frac{AB}{AC}=\frac8{10}=\frac45.
\]

По условию луч $AF$ противоположен лучу $AB$, следовательно,
\[
\angle FAC+\angle BAC=180^\circ.
\]
Значит,
\[
\angle FAC=180^\circ-\angle BAC.
\]

Используем формулы:
\[
\sin(180^\circ-\alpha)=\sin\alpha,
\qquad
\cos(180^\circ-\alpha)=-\cos\alpha.
\]
Тогда
\[
\sin\angle FAC=\sin\angle BAC=\frac35,
\]
а
\[
\cos\angle FAC=-\cos\angle BAC=-\frac45.
\]

Итак,
\[
\boxed{AC=10,\qquad
\sin\angle FAC=\frac35,\qquad
\cos\angle FAC=-\frac45}.
\]

\subsection*{Проверка}
Угол $\angle BAC$ острый, поэтому его синус и косинус положительны.

Угол $\angle FAC$ смежный с острым углом, следовательно, он тупой. У тупого угла синус положителен, а косинус отрицателен. Полученные знаки
\[
\sin\angle FAC>0,\qquad \cos\angle FAC<0
\]
этому соответствуют.

Также проверим основное тригонометрическое тождество:
\[
\left(\frac35\right)^2+\left(-\frac45\right)^2
=\frac9{25}+\frac{16}{25}
=1.
\]

\subsection*{Итог}
\textcolor{teal!60!black}{Итог:}
\[
AC=10,\qquad
\sin\angle FAC=\frac35,\qquad
\cos\angle FAC=-\frac45.
\]

\subsection*{Задача на закрепление}
В прямоугольном треугольнике $PQR$ угол $Q$ равен $90^\circ$, $PQ=12$, $QR=5$. Луч $PS$ является продолжением луча $PQ$ в противоположную сторону. Найдите длину $PR$, а также $\sin\angle SPR$ и $\cos\angle SPR$.

\newpage
\phantomsection\label{task:5d}
\section*{Задание 5г}

\subsection*{Текст задания}
Упростите выражение и укажите ОДЗ:
\[
\frac{5x-3}{3-5x}.
\]

\subsection*{Ваше решение}
Упрощение выполнено так:
\[
\frac{5x-3}{3-5x}
=
\frac{5x-3}{-(5x-3)}
=
-1.
\]
Далее для ОДЗ записано:
\[
3-5x\ne0,
\]
после чего получено
\[
-5x\ne3,
\qquad
x\ne-\frac35.
\]

\subsection*{Ваш ответ}
\[
-1,\qquad x\ne-\frac35.
\]

\subsection*{Ошибка}
\textcolor{red}{Ошибка:} при решении условия для знаменателя потерян знак минус.

\subsection*{Где именно возникает ошибка}
Последний правильный шаг:
\[
3-5x\ne0.
\]

Первый неправильный шаг:
\[
-5x\ne3.
\]

Чтобы убрать число $3$ из левой части, нужно вычесть $3$ из обеих частей неравенства:
\[
3-5x\ne0
\quad\Longrightarrow\quad
-5x\ne-3.
\]
Именно знак перед числом $3$ был потерян.

\subsection*{Почему этот шаг неверен}
Условие существования дроби определяется знаменателем:
\[
3-5x\ne0.
\]
При переносе слагаемого через знак равенства или неравенства его знак меняется. Поэтому из
\[
3-5x\ne0
\]
следует
\[
-5x\ne-3,
\]
а не $-5x\ne3$.

После деления обеих частей на $-5$ получается
\[
x\ne\frac35.
\]
Знак снова становится положительным, потому что
\[
\frac{-3}{-5}=\frac35.
\]

Само упрощение выражения выполнено верно:
\[
3-5x=-(5x-3),
\]
поэтому при допустимых значениях переменной дробь действительно равна $-1$.

Важно сохранить ограничение $x\ne\frac35$: после сокращения исходный знаменатель уже не виден, но запрещённое значение из исходного выражения всё равно остаётся запрещённым.

\subsection*{Ключевая идея}
\textcolor{blue!60!black}{Ключевая идея:} ОДЗ всегда находят по исходному знаменателю до сокращений, а каждый перенос выполняют с обязательной сменой знака переносимого слагаемого.

\subsection*{Исправление от последнего правильного шага}
Начинаем с верной записи:
\[
3-5x\ne0.
\]
Вычитаем $3$ из обеих частей:
\[
-5x\ne-3.
\]
Делим на $-5$:
\[
x\ne\frac35.
\]

Следовательно, окончательно:
\[
\frac{5x-3}{3-5x}=-1,
\qquad
x\ne\frac35.
\]

\subsection*{Полное правильное решение}
Сначала укажем ОДЗ. Знаменатель не должен обращаться в ноль:
\[
3-5x\ne0.
\]
Решаем это условие:
\[
-5x\ne-3,
\]
\[
x\ne\frac35.
\]

Теперь упростим выражение. Заметим:
\[
3-5x=-(5x-3).
\]
Поэтому
\[
\frac{5x-3}{3-5x}
=
\frac{5x-3}{-(5x-3)}.
\]
При $x\ne\frac35$ числитель $5x-3$ не равен нулю, поэтому можно сократить общий множитель:
\[
\frac{5x-3}{-(5x-3)}=-1.
\]

Итак,
\[
\boxed{-1,\qquad x\ne\frac35}.
\]

\subsection*{Проверка}
Возьмём допустимое значение, например $x=0$:
\[
\frac{5\cdot0-3}{3-5\cdot0}
=
\frac{-3}{3}
=
-1.
\]
Это согласуется с упрощённым выражением.

Проверим запрещённое значение:
\[
x=\frac35.
\]
Тогда знаменатель исходной дроби:
\[
3-5\cdot\frac35=3-3=0.
\]
Следовательно, именно $x=\frac35$ должно быть исключено из ОДЗ.

\subsection*{Итог}
\textcolor{teal!60!black}{Итог:}
\[
\frac{5x-3}{3-5x}=-1
\quad\text{при}\quad
x\ne\frac35.
\]

\subsection*{Задача на закрепление}
Упростите выражение и укажите ОДЗ:
\[
\frac{7x-4}{4-7x}.
\]

\vfill
\begin{center}
\small Лёвин Артём Александрович эксклюзивно для Софии Хоменко
\end{center}

\end{document}
